\chapter{Reading results}\label{ch:results}

\section{The Results panel}

Each case run gets a tab in the \gui{Results} panel, titled with the case's title. The
last tab, \gui{Log}, holds the full text of every result, the validation messages and
anything else worth keeping, all selectable and copyable. A result tab has three parts:

\begin{description}
  \item[Headline] the main answer in one or two lines, such as \emph{U\_th = 0.9348 pu <
    -3.0196 deg   Z\_th,1 = 0.1645 + 0.0677i pu}.
  \item[Warnings] in yellow banners: anything you should know, for example an element left
    out, a zero sequence that could not be built, an override that names something gone,
    or a load flow that did not converge. A case that could not run shows only its reason.
  \item[Tables] the numbers, each with its title and sometimes a note below it.
\end{description}

\begin{figure}[H]
  \centering\shot[0.86\linewidth]{results-loadflow.png}
  \caption{The first result tab of the load flow example (top part).}
\end{figure}

\section{Number formats}

\app\ prints numbers the way MATLAB's \code{format short} does, so they can be compared
with a MATLAB solution digit by digit and pasted into a report:

\begin{itemize}
  \item four decimals in fixed point: \code{0.9832}, \code{22.1658};
  \item scientific notation when a value is very large or very small: \code{3.6037e+03},
    \code{1.2830e-05};
  \item impedances, admittances and complex powers in rectangular form,
    \code{0.1645 + 0.0677i}; voltages and currents in polar form with the angle in
    degrees, \code{0.9952 pu < -0.5345 deg};
  \item matrices with one common scale, as MATLAB prints them. When the largest entry
    leaves the fixed-point window, the title carries the factor (\code{1.0e+02 *}) and the
    entries are divided by it.
\end{itemize}

Units are always given: in the column header, after the value, or in the table's note.

\section{Copying}

\begin{itemize}
  \item Select cells and press \keys{\ctrl + C} to copy them as tab-separated text, ready
    for a spreadsheet.
  \item Right-click a table for \gui{Copy}, the selection, or \gui{Copy table}, all of it
    with its header row.
  \item The headline can be selected with the mouse and copied.
\end{itemize}

\subsection{Copy as MATLAB}\label{sec:copy-matlab}

A table that prints a matrix (an admittance matrix, the load-flow Jacobian, the S1a
matrix $F$, the $Z_\text{ph}$ and $Z_s$ of an unbalanced load) offers \gui{Copy as
MATLAB} in its context menu. It copies the matrix as MATLAB source with every digit kept,
not the four decimals on screen. The two-port matrix of S1, for example, copies as:

\begin{terminal}
F = [1.0002045432793645 + 0.0004317301765735776i, 0.004880874748045769 + 0.27917704044800307i; ...
     0.00369999999509357 - 0.0017535002966172987i, 1.0003032467164057 + 0.000592414161302107i];
\end{terminal}

Paste it into MATLAB and compare it with your own matrix to the last digit.

\section{Annotating the diagram}\label{sec:annotations}

Most studies can annotate the diagram with their result. After a run, the selected result
tab's annotations are drawn next to the elements they describe:

\begin{itemize}
  \item bus annotations above the bus, in green;
  \item branch annotations below the branch, beside it for a vertical branch, in brown;
  \item element annotations, such as a generator's $P$ and $Q$, to the right of the
    symbol;
  \item values that call for attention, such as a faulted bus or a voltage outside
    0.95--1.05~pu, in the warning colour.
\end{itemize}

Click another result tab and the diagram shows that result instead. \menu{View > Show
results on the diagram} (\keys{\ctrl + R}) hides or shows the annotations. They follow
their element when it moves, and they disappear when the model changes, because they no
longer describe it. Each study's section in \cref{ch:studies} says what it draws. The
model, sequence model, admittance matrix, S1a and S1b results are not drawn on the
diagram.

\begin{figure}[H]
  \centering\shot[0.92\linewidth]{diagram-s3c.png}
  \caption{A single line-to-ground fault at bus 7 of \file{grid\_industry.psn} (case
    S3c). The faulted bus shows the fault current and its three phase voltages, the other
    buses their phase voltages during the fault, and the branches their largest current.}
\end{figure}

\section{Colouring the diagram}\label{sec:colouring}

\menu{View > Colour by} colours the diagram, with a legend in the top-right corner of the
canvas:

\begin{description}
  \item[Nothing] every element in black.
  \item[Voltage level] one colour per base voltage, highest first. Buses, the elements that
    sit in one voltage zone and the wires that land on each bus take their zone's colour.
    Transformers, which join two zones, stay black. The colours follow edits: a new
    element is coloured as soon as it is connected.
  \item[Voltage result] every bus coloured by the band its voltage magnitude falls in, for
    the result tab shown: below 0.90, 0.90--0.95, 0.95--0.98, 0.98--1.02, 1.02--1.05 and
    above 1.05~pu. During a fault, the lowest phase counts. A result without bus
    voltages, such as an admittance matrix, leaves the buses black and says so in the
    status bar.
\end{description}

\begin{figure}[H]
  \centering
  \begin{minipage}[t]{0.49\linewidth}
    \centering\shot[\linewidth]{colour-level.png}
  \end{minipage}\hfill
  \begin{minipage}[t]{0.49\linewidth}
    \centering\shot[\linewidth]{colour-voltage.png}
  \end{minipage}
  \caption{Left: the S2 grid with the S1 industry coloured by voltage level (70, 40, 12
    and 6~kV). Right: the load flow coloured by bus voltage. Bus~8, at 0.9661~pu, falls in
    the 0.95--0.98 band.}
\end{figure}

\section{Validation messages}

Before every run, \app\ validates the network. Errors stop the run and are listed in a
message box and in the \gui{Log}. A component that is not wired does not stop the run:
the studies leave it out, and every result says so in a warning, \emph{``LD9 is not
connected to a bus, so the studies ignore it.''} \Cref{ch:troubleshooting} lists the
messages with what to do.
